\section*{Capítulo \ref{ch:b3:sensory-systems} --- Sistemas sensoriais}

\begin{solution}{exo:b3:sensory-systems:1}
Pelo fio, e não pela mensagem: cada nervo sensorial é uma \omterm{def:b3:sensory-systems:principles}{linha rotulada}
que termina em sua própria região do encéfalo, e o que quer que chegue
pelo nervo óptico é lido como luz --- razão pela qual uma pressão sobre
o globo ocular produz um lampejo e um golpe no ouvido, um zumbido.
\end{solution}

\begin{solution}{exo:b3:sensory-systems:2}
No escuro, o GMP cíclico mantém abertos canais de cátions, e uma
``corrente de escuro'' constante para dentro mantém a célula
despolarizada e liberando glutamato. A luz ativa a \omterm{def:b3:sensory-systems:retina}{rodopsina}, a
transducina e a fosfodiesterase, que destrói o cGMP; os canais fecham, a
corrente para dentro cessa e a célula hiperpolariza e libera menos --- a
luz é sinalizada por uma diminuição.
\end{solution}

\begin{solution}{exo:b3:sensory-systems:3}
$I = I_{0}10^{L/10}$: \qty{20}{dB}, \qty{1e-10}{W/m^2}; \qty{60}{dB},
\qty{1e-6}{W/m^2}; \qty{110}{dB}, \qty{0.1}{W/m^2}. Do mais alto ao mais
baixo: $10^{9}$.
\end{solution}

\begin{solution}{exo:b3:sensory-systems:4}
Doce, umami e amargo por receptores acoplados a proteína~G (um de doce,
um de umami, cerca de vinte e cinco de amargo); salgado pelo canal de
sódio ENaC; ácido por um canal de prótons. A maior parte do ``sabor'' é
o cheiro de moléculas voláteis que chegam ao nariz a partir da boca; com
o nariz entupido restam apenas os cinco gostos e a comida parece sem
graça.
\end{solution}

\begin{solution}{exo:b3:sensory-systems:5}
Fechner: $S = (1/0.02)\ln(10\,000/100) = 50\ln 100 \approx 230$ passos.
Contando: $1.02^{n} = 100$ dá $n = \ln 100/\ln 1.02 = 233$.
\end{solution}

\begin{solution}{exo:b3:sensory-systems:6}
$P = 1 - e^{-a}\sum_{k=0}^{5}a^{k}/k!$: $a = 2$: $0.017$; $a = 6$:
$0.55$; $a = 12$: $0.98$. Metade dos lampejos é vista em $a \approx 5.7$;
com \qty{10}{\%} de absorção o lampejo precisa entregar cerca de $57$
fótons na córnea.
\end{solution}

\begin{solution}{exo:b3:sensory-systems:7}
$55/3\times 1.3 \approx 24$; em decibéis de pressão, $20\log_{10}24
\approx \qty{28}{dB}$. Ossículos fundidos não conseguem transmitir a
vibração, de modo que o som só chega à \omterm{def:b3:sensory-systems:cochlea}{cóclea} por condução pelo crânio,
sem o ganho do ouvido médio: uma perda condutiva de uns
\qtyrange{30}{50}{dB}.
\end{solution}

\begin{solution}{exo:b3:sensory-systems:8}
O canal da \omterm{def:b3:sensory-systems:cochlea}{célula ciliada} é aberto diretamente pela \omterm{def:b3:sensory-systems:cochlea}{ligação de ponta} que
o puxa --- sem mensageiro, sem enzima ---, de modo que responde em
microssegundos; o fotorreceptor amplifica um fóton por duas etapas
enzimáticas, que levam dezenas de milissegundos. A audição ganha a
precisão de temporização necessária para localizar sons por atrasos
interaurais e para acompanhar a fase; a visão ganha a sensibilidade para
contar fótons isolados, e paga em rapidez.
\end{solution}

\begin{solution}{exo:b3:sensory-systems:9}
$H(0.05) = -0.05\ln 0.05 - 0.95\ln 0.95 = 0.150 + 0.049 = 0.199$.
Humano: $\ln\binom{400}{20} \approx 400\times 0.199 = 80$, cerca de
$10^{34}$ padrões. Camundongo: $\ln\binom{1100}{55} \approx 1100\times 0.199
= 219$, cerca de $10^{95}$.
\end{solution}

\begin{solution}{exo:b3:sensory-systems:10}
Média de $500\times 0.1/40 = 1.25$ eventos espontâneos por janela. $P(k
\le 5) = e^{-1.25}(1 + 1.25 + 0.78 + 0.33 + 0.10 + 0.03) = 0.998$, de
modo que $P(k \ge 6) \approx 2\times 10^{-3}$. Com um limiar de um, um
evento espontâneo ocorreria na maioria das janelas ($1 - e^{-1.25} = 0.71$) e o escuro estaria cheio de lampejos; em seis, os alarmes falsos
vêm uma vez em quinhentas janelas. O limiar é posto onde o ruído dos
\omterm{def:b3:sensory-systems:retina}{bastonetes}, e não sua sensibilidade, manda.
\end{solution}

\begin{solution}{exo:b3:sensory-systems:11}
Maior atraso: $0.2/340 = \qty{590}{\micro s}$ (som vindo de um lado).
Perto da linha média $\Delta t \approx (d/c)\sin\theta$, de modo que
$\qty{10}{\micro
s}$ correspondem a $\sin\theta = 0.017$, cerca de
\qty{1}{\degree}. Acima de \qty{4}{kHz} o período
(\qty{250}{\micro s}) é mais curto do que os neurônios conseguem travar,
e um atraso de várias centenas de microssegundos abrange mais de um
ciclo, tornando a fase ambígua; o encéfalo usa então a diferença de
nível entre os ouvidos e o código de lugar.
\end{solution}

\begin{solution}{exo:b3:sensory-systems:12}
Longe da borda, $r = I(1 - 2w)$: $0.6$ do lado escuro, $1.2$ do lado
claro. No último receptor escuro ($I = 1$, vizinhos $1$ e $2$): $r
= 1 - 0.2\times 3 = 0.4$; no primeiro claro: $r = 2 - 0.6 = 1.4$. O
degrau é exagerado num subdisparo e num sobredisparo --- as bandas de
Mach. A \omterm{def:b3:sensory-systems:retina}{retina} ganha contraste e bordas, e uma faixa dinâmica menor para
transmitir; perde fidelidade aos níveis absolutos e produz ilusões de
brilho.
\end{solution}

\begin{solution}{pb:b3:sensory-systems:1}
\textbf{1.} $0.01/3.6\times 10^{-19} = 2.8\times 10^{16}$ fótons por
segundo.
\textbf{2.} Área da pupila: $\pi(3.5\times 10^{-3})^{2} = 3.8\times 10^{-5}$ m$^{2}$. A \qty{1}{km}: fração $3.8\times 10^{-5}/1.26\times 10^{7} =
3\times 10^{-12}$, $8.5\times 10^{4}$ fótons por segundo. A
\qty{10}{km}: $850$ por segundo.
\textbf{3.} Em \qty{0.1}{s} com \qty{10}{\%} absorvidos: $850$ a
\qty{1}{km}, $8.5$ a \qty{10}{km} --- os dois acima de seis; a vela é
vista a dez quilômetros (na maior parte das vezes).
\textbf{4.} Seis absorvidos por janela exigem $600$ fótons por segundo
entrando na pupila: $4\pi d^{2} = 3.8\times 10^{-5}\times 2.8\times 10^{16}/600
= 1.8\times 10^{9}$ m$^{2}$, $d \approx \qty{12}{km}$. Ali,
$P(k \ge
6 \mid a = 6) = 0.55$: a vela é vista em cerca de metade das
janelas.
\textbf{5.} $e^{-6} = 0.0025$. No limiar a contagem flutua de janela em
janela --- $6 \pm 2.4$ ---, de modo que a fonte parece ir e vir: a
cintilação de uma estrela fraca (à qual a atmosfera acrescenta a sua).
\textbf{6.} A flutuação do fundo, $\sqrt{10^{4}} = 100$ fótons por
janela, engole os $8.5$ da vela: a detecção exige que o sinal supere o
ruído do fundo, e não o limiar do olho no escuro.
\textbf{7.} $\qty{1}{pA}/200 = \qty{5}{fA}$ por canal; $5\times
10^{-15}/1.6\times 10^{-19} = 3\times 10^{4}$ íons por segundo.
\textbf{8.} Um trinta avos; cerca de trinta fótons simultâneos fecham
todos os canais e saturam o \omterm{def:b3:sensory-systems:retina}{bastonete}.
\textbf{9.} $10^{-12}\times 0.2 = 2\times 10^{-13}$ C, $1.3\times 10^{6}$
cátions: um fóton controla mais de um milhão de cargas.
\textbf{10.} O \omterm{def:b3:sensory-systems:retina}{bastonete} satura em milissegundos, com todos os canais
fechados, e sua \omterm{def:b3:sensory-systems:retina}{rodopsina} é branqueada mais depressa do que regenerada;
os \omterm{def:b3:sensory-systems:retina}{cones}, de ganho menor e desligamento mais rápido, continuam a
responder e deslocam sua faixa ao longo de seis décadas de luz.
\textbf{11.} Ganho alto exige integração longa (cada etapa leva tempo e a
resposta é prolongada para acumular), de modo que o \omterm{def:b3:sensory-systems:retina}{bastonete} é lento e
sensível; a amplificação menor do \omterm{def:b3:sensory-systems:retina}{cone} e seu desligamento mais rápido
dão uma resposta em \qty{20}{ms}, rápida o bastante para o movimento, ao
custo de precisar de cem fótons.
\textbf{12.} $10^{8}/40 = 2.5\times 10^{6}$ isomerizações espontâneas por
segundo em toda a \omterm{def:b3:sensory-systems:retina}{retina} --- mas apenas $1.25$ por janela em qualquer
retalho de $500$ \omterm{def:b3:sensory-systems:retina}{bastonetes}, muito abaixo do limiar de seis, de modo que
o ruído é rejeitado retalho por retalho; o limiar e o agrupamento
existem precisamente para isso.
\textbf{13.} \qty{0.1}{W/m^2}; sobre \qty{5.5e-5}{m^2},
\qty{5.5e-6}{W}.
\textbf{14.} $5.5\times 10^{-6}\times 7200 = \qty{0.04}{J}$ --- cerca de
quatrocentas gotas de chuva ao longo de duas horas.
\textbf{15.} $10^{-12}\times 5.5\times 10^{-5} = 5.5\times 10^{-17}$ W;
em \qty{10}{ms}, $5.5\times 10^{-19}$ J --- cerca da energia de um ou
dois fótons visíveis.
\textbf{16.} $0.3/5000 = 6\times 10^{-5}$ rad, $0.003^{\circ}$; a
deflexão equivale a uns três átomos de hidrogênio.
\textbf{17.} $10^{(110 - 85)/10} = 316$: duas horas a \qty{110}{dB}
entregam a energia de $630$ horas a \qty{85}{dB} --- quase um mês de
dias de trabalho.
\textbf{18.} Cerca de $120$ passos apenas perceptíveis, um por \omterm{def:b3:sensory-systems:cochlea}{decibel}.
\textbf{19.} $2^{400} = 10^{400\log_{10}2} = 10^{120}$.
\textbf{20.} $\ln\binom{400}{20} \approx 400\,H(0.05) \approx 80$, menos
uma correção de Stirling de cerca de $2$: uns $10^{34}$ padrões.
\textbf{21.} Conte os receptores presentes num padrão e ausentes no
outro: $5 +
5 = 10$ de $40$ (uma distância de Hamming); padrões que se
sobrepõem em três quartos acionam quase os mesmos \omterm{def:b3:sensory-systems:chemical}{glomérulos} e são
lidos como quase o mesmo cheiro.
\textbf{22.} $10^{12}/10^{34} = 10^{-22}$ dos padrões. O limite não é o
código: os receptores são ruidosos, muitos têm sintonia parecida, de
modo que suas atividades são correlacionadas, as moléculas reais
produzem apenas um pequeno subconjunto de padrões, e a discriminação e a
memória de padrões do encéfalo são finitas.
\textbf{23.} $\ln\binom{1100}{20} \approx 1100\,H(0.018) = 1100\times
0.091 \approx 100$, cerca de $10^{43}$ --- uns $10^{9}$ vezes o número
humano de padrões de vinte receptores (e muito mais se o camundongo usar
mais receptores por odor).
\textbf{24.} O olfato em geral quase não muda --- o código é redundante e
outros receptores carregam todo odorante ---, mas um odorante que
dependa desse receptor em baixa concentração se torna indetectável ou
muda de caráter: uma anosmia específica, como nas muitas pessoas que não
conseguem sentir o cheiro da androstenona.
\textbf{25.} A vela é vista até cerca de \qty{12}{km}; o som no limiar
entrega $5.5\times 10^{-19}$ J ao tímpano em \qty{10}{ms}, um fóton ou
dois; um nariz humano consegue formar uns $10^{34}$ padrões de vinte
receptores.
\end{solution}
